% !TEX root = ../DoAn.tex
\documentclass[../DoAn.tex]{subfiles}
\begin{document}

\section{Tổng kết}
\label{section:5.1}

Đồ án giải quyết bài toán hỏi đáp trên kho tài liệu kỹ thuật tiếng Việt bằng kiến trúc Retrieval-Augmented Generation, với trọng tâm là tối ưu hai thành phần quyết định chất lượng truy hồi: mô hình embedding và reranker. Xuất phát từ một hệ thống nền tảng (v1), đồ án đo lường và xác định ba điểm nghẽn cụ thể về độ trễ, chi phí và chất lượng, sau đó đề xuất hệ thống v2 với nguyên tắc mỗi thay đổi kỹ thuật phải giải quyết một vấn đề đã được đo lường thay vì tối ưu dựa trên giả định.

Về embedding, đồ án tinh chỉnh mô hình \plaincode{Vietnamese_Embedding_v2} (dựa trên BGE-M3) bằng hàm mất mát MultipleNegativesRankingLoss trên các cặp câu hỏi--positive của miền, kết hợp Matryoshka Representation Learning. Mô hình sau tinh chỉnh nâng Acc@1 trên tập in-domain từ 65.19\% lên $73.21 \pm 0.60$\% (trung bình trên ba seed), đồng thời vượt cả BGE-M3 gốc trên tập tài liệu độc lập, cho thấy việc thích nghi miền có hiệu quả nhất quán chứ không chỉ trên dữ liệu quen thuộc. Nhờ MRL, biểu diễn 512 chiều giữ được phần lớn chất lượng của 1024 chiều, giúp giảm một nửa dung lượng vector index mà không hy sinh đáng kể độ chính xác.

Về reranker, đồ án chưng cất teacher BGE-reranker-v2-m3 (568M tham số) thành các student nhẹ thông qua quy trình hai giai đoạn (contrastive learning rồi chưng cất tri thức) với hàm mất mát ADR-MSE theo hướng rank-based. Student mmarco 117M đạt MRR@10 = 0.8862, bằng đúng teacher 568M, còn phiên bản cắt tỉa từ vựng 35.6M đạt MRR@10 = 0.8844 với dung lượng nhỏ hơn nhiều. Đáng chú ý, reranker 35.6M vượt ViRanker 568M (0.8593, lớn hơn khoảng 16 lần) trên cùng tập đánh giá. Sau lượng tử hoá INT8, mô hình chỉ còn 34.8 MB và chạy trên CPU nhanh hơn teacher khoảng 10 lần, loại bỏ hoàn toàn yêu cầu GPU khỏi tầng xếp hạng lại.

Về routing, đồ án thay GPT router bằng một bộ phân loại ý định cục bộ dựa trên Sentence-BERT, loại bỏ phụ thuộc vào API ngoài, giảm độ trễ và chi phí trên mỗi truy vấn đồng thời giữ dữ liệu trong hạ tầng nội bộ.

\section{Đóng góp chính}
\label{section:5.2}

Các đóng góp chính của đồ án gồm:

\begin{itemize}
\item \textbf{Quy trình xây dựng dữ liệu huấn luyện dùng chung cho embedding và reranker.} Quy trình dùng mô hình ngôn ngữ làm giám khảo xác định positive, lọc khả năng trả lời độc lập, rồi khai thác hard negative từ kết quả truy hồi cho tầng xếp hạng. Cùng một nguồn dữ liệu được tái sử dụng cho cả hai tầng với tiêu chí lọc riêng.
\item \textbf{Tinh chỉnh embedding theo miền kết hợp MRL.} Tinh chỉnh bằng MultipleNegativesRankingLoss trên các cặp câu hỏi--positive của miền, kết hợp MRL để một mô hình duy nhất cho biểu diễn dùng được ở ba số chiều; nhờ đó embedding 512 chiều giữ được phần lớn chất lượng của 1024 chiều với nửa dung lượng index.
\item \textbf{Nén reranker bằng chưng cất tri thức kết hợp cắt tỉa từ vựng.} Chưng cất teacher BGE-reranker-v2-m3 qua hai giai đoạn rồi cắt tỉa từ vựng, rồi lượng tử hoá INT8, tạo ra reranker 34.8 MB đạt chất lượng cao hơn các baseline tiếng Việt lớn hơn nhiều lần và chạy được hoàn toàn trên CPU.
\item \textbf{Hệ thống RAG tối ưu đầu cuối.} Tích hợp embedding tinh chỉnh 512 chiều, reranker chưng cất và bộ phân loại ý định cục bộ thay GPT router, giảm độ trễ, dung lượng và chi phí suy luận so với hệ thống v1 trong khi vẫn giữ chất lượng truy hồi và xếp hạng.
\end{itemize}

\section{Hạn chế}
\label{section:5.3}

Bên cạnh các kết quả đạt được, đồ án còn một số hạn chế.

Thứ nhất, mức cải thiện của tinh chỉnh embedding rõ rệt trên tập in-domain nhưng khiêm tốn trên tập tài liệu độc lập (khoảng một điểm phần trăm Acc@1), do base model đã được huấn luyện contrastive ở quy mô lớn và lượng dữ liệu miền còn nhỏ. Bi-encoder cũng có trần kiến trúc khi phải nén toàn bộ ngữ nghĩa vào một vector đơn, khiến việc phân biệt các hard negative khác nhau ở mức chi tiết factual (số liệu, thực thể) trở nên khó khăn; phần này được bù đắp bởi reranker nhưng chưa được giải quyết triệt để ở tầng embedding.

Thứ hai, cả dữ liệu huấn luyện lẫn dữ liệu đánh giá đều có một phần được sinh tự động qua mô hình ngôn ngữ và xác nhận thủ công bởi chính tác giả. Với 50 tài liệu của tập D2 --- vốn không có bộ câu hỏi--đáp án tham chiếu sẵn --- toàn bộ câu hỏi và đáp án gold dùng cho đánh giá reranker, embedding và MCQ đầu cuối (mục~\ref{subsection:4.1.2}, \ref{subsection:4.2.5}) do tác giả tự biên soạn với sự hỗ trợ của GPT, sau đó tự kiểm tra thủ công. Người kiểm tra không độc lập với người đánh giá hệ thống, nên không loại trừ hoàn toàn khả năng thiên lệch trong cách diễn đạt câu hỏi. Thêm vào đó, bộ 390 câu MCQ vừa được dùng để lựa chọn cấu hình embedding/reranker vừa được dùng lại để báo cáo kết quả đầu cuối, nên các chỉ số tuyệt đối có thể lạc quan hơn thực tế do hiệu ứng lựa chọn cấu hình trên chính tập báo cáo. Để định lượng ảnh hưởng này, một mẫu đối chiếu độc lập gồm 71 câu (chưa từng dùng ở bất kỳ bước lựa chọn cấu hình nào) đã được đánh giá riêng ở mục~\ref{subsection:4.2.7}: kết quả cho thấy độ chính xác giảm khoảng 4--5 điểm phần trăm so với bộ 390 câu, xác nhận mức lạc quan như dự đoán, nhưng kết luận so sánh giữa v1 và v2 --- độ chính xác tương đương, tốc độ nhanh hơn đáng kể, router cục bộ chính xác hơn GPT router --- vẫn tái lập nhất quán trên dữ liệu độc lập này.

Thứ ba, cắt tỉa từ vựng giảm dung lượng mô hình nhưng bản thân nó không giảm đáng kể độ trễ của Transformer encoder, do số tầng và kích thước ẩn được giữ nguyên; phần tăng tốc chủ yếu đến từ chưng cất (student nhỏ hơn teacher) và lượng tử hoá INT8. Lượng tử hoá được dùng ở đây là dạng động (chỉ nén trọng số), chưa khai thác hết mức tăng tốc mà lượng tử hoá tĩnh có thể mang lại.

Thứ tư, hệ thống được đánh giá trên một miền tài liệu kỹ thuật cụ thể. Khả năng tổng quát hóa sang các miền khác hoặc các loại tài liệu khác cần được kiểm chứng thêm.

\section{Hướng phát triển}
\label{section:5.4}

Từ các hạn chế trên, đồ án đề xuất ba hướng phát triển tiếp theo.

Thứ nhất, về dữ liệu và chất lượng truy hồi, mẫu đối chiếu độc lập ở mục~\ref{subsection:4.2.7} là một bước đầu để kiểm chứng độ tin cậy của đánh giá; hướng xa hơn là xây dựng một bộ kiểm thử được một người chú thích độc lập (không tham gia phát triển hệ thống) xác minh hoàn toàn, nhằm loại bỏ khả năng thiên lệch do người đánh giá đồng thời là người biên soạn câu hỏi. Đồng thời, hướng khả dĩ để nâng trần chất lượng ở tầng truy hồi là mở rộng số lượng câu hỏi thật của miền, thay vì tăng độ phức tạp của hàm mất mát.

Thứ hai, về mô hình, lượng tử hoá tĩnh (nén cả activation, cần tập dữ liệu hiệu chuẩn) có thể đạt mức tăng tốc cao hơn lượng tử hoá động đang dùng; ngoài ra có thể khảo sát cắt tỉa tầng kèm chưng cất, cũng như thử nghiệm các kiến trúc reranker khác để tìm điểm cân bằng tốt hơn giữa chất lượng và tốc độ.

Thứ ba, về khả năng tổng quát hóa và triển khai thực tế, việc đánh giá hệ thống trên nhiều miền tài liệu khác nhau và thử nghiệm trong môi trường vận hành thực tế sẽ giúp kiểm chứng tính ứng dụng của các giải pháp đề xuất ngoài phạm vi miền tài liệu kỹ thuật hiện tại.

\end{document}